\section{Related Work}\label{sec:related}

\subsection{Determinants of higher-education participation}

The literature on Indian higher education groups the determinants of participation into demand-side and supply-side factors. Demand-side accounts emphasise household income, secondary completion, aspirations, caste and gender norms, and the expected return to a degree \cite{sharma2026,ghosh2021}. Supply-side accounts emphasise the number and location of institutions, the availability of teachers, programme mix, distance education and affirmative-action capacity. Recent State-level work reaches similar conclusions. Sharma et al.~\cite{sharma2026} show that private provision narrows regional GER gaps but not socio-economic gaps; Kadlak and Behera~\cite{kadlak2026} argue that the aggregate GER masks district-level access in Maharashtra; Mitra and Ghara~\cite{mitra2019} document district-level variation within West Bengal; and Baruah and Goswami~\cite{baruah2024} evaluate the effect of the Rashtriya Uchchatar Shiksha Abhiyan on GER. Ghosh and Kundu~\cite{ghosh2021} find that women's participation varies sharply across major States. Outside India, Akinola and Mbonigaba~\cite{akinola2020} use panel data for Sub-Saharan Africa and find that income and secondary completion dominate, while Prakasam et al.~\cite{prakasam2019} show that discipline-level enrolment responds differently to labour-market signals.

The policy-visible supply indicators used here are college density (colleges per lakh of the 18--23 population), PTR and the Gender Parity Index (GPI, female GER divided by male GER). Each is an imperfect proxy. Density measures nominal access, not capacity, quality or programme fit. PTR is an input measure that can worsen mechanically when enrolment grows faster than teacher hiring. GPI summarises the gender gap in aggregate participation but not field-of-study segregation or completion. These limitations motivate a multivariate and explicitly conditional analysis rather than a bivariate ranking of correlates.

\subsection{Machine learning and policy analysis}

Machine learning is increasingly used in policy analysis, but its comparative advantage is prediction rather than causal identification \cite{mullainathan2017,athey2019}. A flexible model can improve out-of-sample forecasts while remaining uninformative about the effect of an intervention. Conversely, a transparent linear model can be preferable for communication and extrapolation even when a flexible learner has a slightly lower cross-validated error. This tension is central here: the polynomial model has the best LOOCV error, but it is not used for projection because it becomes non-monotone outside the observed density range. The linear model is retained because it is interpretable, monotone and defensible at the boundary of the data. The same distinction appears in learning-analytics work on equitable enrolment, where accurate prediction of who enrols does not by itself identify which allocation rule is fair \cite{alenezi2026}.

The statistical toolkit is standard. Ridge regression controls variance in the presence of collinear predictors by adding an $L_2$ penalty \cite{hoerl1970}. Leave-one-out cross-validation provides an almost unbiased estimate of prediction error for small samples \cite{stone1974}, while repeated $k$-fold cross-validation checks whether a ranking is an artefact of one partition. Heteroskedasticity-robust covariance estimators \cite{white1980,mackinnon1985}, the bootstrap \cite{efron1993}, Cook's distance \cite{cook1977} and the classical specification tests \cite{shapiro1965,breusch1979,ramsey1969} are used to interrogate the fitted model rather than to decorate it. Panel fixed effects and first differences follow the standard within-estimator logic \cite{angrist2009}, with cluster-robust inference \cite{cameron2015}.

\subsection{Gap addressed}

Three gaps recur in applied work on GER. First, many analyses are bivariate or use national time series, so they cannot say how much independent information each indicator carries. Second, out-of-sample testing, influence analysis and panel checks are often omitted, which makes a fragile cross-sectional correlation look like a structural relationship. Third, policy discussions frequently jump from a correlation to a 2035 projection without stating the trend assumption or the uncertainty. This study addresses all three with a small, fully documented dataset and a protocol that a reader can re-run.
